\documentclass{article}
\usepackage{graphicx} % Required for inserting images

\title{NLP exam project}
\author{Nada Bou Kanaan}
\date{May 2026}



\begin{document}

\maketitle

\section{Goal}

Implement and evaluate different prompting strategies for Text-to-Cypher generation on the Neo4j \texttt{recommendations} graph, using Qwen as the base model. The project also includes a supervised Cypher difficulty classification baseline in order to strengthen the NLP component of the homework.

\section{Problem}

Text-to-Cypher systems must translate natural-language questions into executable graph queries while respecting schema constraints. In practice, this is difficult because small mistakes in labels, properties, filters, aggregation, or return fields can make a query incorrect even when it looks plausible.

Recent work highlights three useful directions. Chauhan et al.~\cite{chauhan2025mind} emphasize the need for realistic benchmark pairs of natural-language questions and validated Cypher queries. Hornsteiner et al.~\cite{hornsteiner2024realtime} stress the importance of strong schema grounding. Soleymani et al.~\cite{soleymani2025tasksplitting} propose decomposition as a way to make Cypher generation easier. Mandilara et al.~\cite{mandilara2025decoding} focus on evaluation and schema-consistency issues.

The main question of this homework is therefore whether more complex prompting strategies actually improve Cypher generation compared with a simpler schema-aware baseline.

\section{Data and Setup}

Two types of data were used in this project. For the Text-to-Cypher generation experiments, I used a reduced benchmark of 30 natural-language questions paired with gold Cypher queries.
The subset preserves easy, medium, and hard questions, together with different query types such as retrieval, aggregation, verbose questions, and more complex graph patterns.

All experiments used the full schema of the Neo4j \texttt{recommendations} database. This was necessary because many benchmark questions require properties such as \texttt{runtime}, \texttt{budget}, \texttt{revenue}, \texttt{languages}, \texttt{imdbVotes}, \texttt{countries}, and \texttt{bio}.

The generation model used in the experiments was Qwen.

For the Cypher difficulty classification experiment, I used a combined dataset built from two sources: the Neo4j Text2Cypher dataset  \cite{ozsoy2025text2cypher},, which I first processed and labeled according to the target difficulty classes, and a manually curated dataset of 2,500 queries

\section{Prompting Strategies}

\subsection{Prompt 1: direct schema-aware prompting}

In Prompt 1, the model receives the full schema, a few-shot set of examples, and the question, and then generates the Cypher query directly in one step.

This strategy is closest to the schema-aware setting discussed by Hornsteiner et al.~\cite{hornsteiner2024realtime}. It is also compatible with the benchmark philosophy of Chauhan et al.~\cite{chauhan2025mind}, where the model is expected to generate correct Cypher under explicit schema constraints.

Strengths: simple, coherent, and fully grounded in the schema. Weaknesses: it may still omit important return fields or mis-handle more subtle semantic constraints.

\subsection{Prompt 2: task splitting}

In Prompt 2, the model first identifies the relevant schema elements and then generates the final Cypher query using those selected elements.

This strategy is directly inspired by Soleymani et al.~\cite{soleymani2025tasksplitting}, who argue that decomposition can help align the question with the graph structure by breaking the generation into smaller subtasks.

Strengths: more controlled schema selection. Weaknesses: mistakes in the first step can propagate into the final query.

\subsection{Prompt 3: iterative refinement}

In Prompt 3, the model first generates a Cypher query, then verifies it internally, and finally attempts to correct it.

This strategy is motivated by the validation-oriented perspective discussed by Mandilara et al.~\cite{mandilara2025decoding}, where Text-to-Cypher should be evaluated not only as generation, but also in terms of schema consistency and semantic correctness.

Strengths: allows self-correction. Weaknesses: the same model acts as both generator and verifier, so deeper errors may not be repaired.

\section{Evaluation}

At first, exact match between generated and gold Cypher was considered as a metric. However, this is too strict, because two queries may differ in aliases or formatting while still being functionally equivalent.

For this reason, the final evaluation was execution-based. For each example:
\begin{enumerate}
    \item the gold query was executed in Neo4j;
    \item the predicted query was executed in Neo4j;
    \item the prediction was counted as correct if it produced the same output as the gold query.
\end{enumerate}

This criterion is more suitable for Text-to-Cypher because it measures functional correctness rather than textual similarity.

The notebook result files still store the generated Cypher strings together with parsing information and exact-match indicators. However, the final performance numbers reported in this homework are based on manual execution in Neo4j and output comparison with the gold queries, rather than on exact string matching alone.

\section{Results}

Using execution-based evaluation on the 30-query dataset, the results were:
\begin{itemize}
    \item Prompt 1: 21/30 correct, equal to 70.0\%
    \item Prompt 2: 16/30 correct, equal to 53.3\%
    \item Prompt 3: 16/30 correct, equal to 53.3\%
\end{itemize}

These results show that Prompt 1 was clearly the strongest of the three strategies. The direct schema-aware approach outperformed both task splitting and iterative refinement.

\section{Discussion}

The experiments show that increasing prompting complexity does not automatically improve performance. Prompt 1 worked best because it preserved the full semantic context of the question and schema in a single generation step. This is consistent with the emphasis on schema grounding found in Hornsteiner et al.~\cite{hornsteiner2024realtime}.

Prompt 2 did not improve over Prompt 1, even though it was inspired by the task-splitting strategy of Soleymani et al.~\cite{soleymani2025tasksplitting}. In this homework setting, decomposition introduced an additional source of fragility, since an incomplete first step could negatively affect the final query.

Prompt 3 also failed to improve over Prompt 1. This suggests that self-correction alone is not sufficient when the same model is responsible for both initial generation and validation. In practice, the iterative-refinement setup did not reliably repair semantic mistakes.

\section{Cypher Difficulty Classification Baseline}

To complement the generation experiments, I also implemented a supervised neural baseline for Cypher difficulty classification. The task is to predict the structural level of a Cypher query, such as node-level, one-hop, two-hop, three-hop, or hard-level, from the query text alone.

The baseline was trained on a combined dataset formed by the Neo4j Text2Cypher dataset and an additional manually curated dataset of 2{,}500 entries. The manually curated set is balanced across difficulty levels and was created to provide cleaner labels and better coverage of difficult cases. Each query was converted into a TF-IDF representation using unigrams and bigrams, and these vectors were used as input to a multilayer perceptron with two hidden layers, ReLU activations, and dropout. The network was trained with cross-entropy loss.

Results:
\begin{itemize}
    \item On the combined dataset: Accuracy = 0.9727, F1 = 0.9529
    \item On the curated hard-level subset: Accuracy = 0.972, Precision = 1.0000
\end{itemize}

This experiment shows that Cypher difficulty can be predicted very effectively even with a relatively simple from-scratch neural architecture. It therefore complements the main generation task by showing that Cypher queries contain strong textual signals of structural complexity.

\section{Proposed Prompt 4}

After evaluating the three prompts, I proposed a new hybrid Prompt 4. Its design was based on the strengths and weaknesses of the previous strategies.

From Prompt 1, it inherits strong schema grounding. From Prompt 2, it inherits the idea that the model should internally identify the minimal relevant schema elements before writing the query. From Prompt 3, it inherits the idea of internal verification before producing the final answer.

At the same time, Prompt 4 avoids the main weaknesses of Prompt 2 and Prompt 3: it does not expose explicit intermediate outputs, and it does not rely on a long multi-step correction process. Instead, it keeps planning and verification internal and produces only one final Cypher query.

Prompt 4 was evaluated on the full 30-query dataset. It produced 4 incorrect outputs and therefore correctly solved 26 out of 30 queries, achieving 86.7\% accuracy. Compared with the previous prompting strategies, this is the best result obtained in the project and suggests that the hybrid design successfully combines the strongest aspects of the earlier prompts.

\section{Conclusion}

This homework shows that execution-based evaluation is more appropriate than exact string matching for Text-to-Cypher generation. Among the three initially tested prompting strategies, direct schema-aware prompting achieved the best result, while task splitting and iterative refinement did not improve performance for Qwen on this dataset.

The additional Cypher difficulty classification experiment showed that the structural complexity of Cypher queries can be modeled effectively by a supervised neural baseline trained from scratch. Finally, based on the strengths and weaknesses of the three prompting strategies, a fourth hybrid prompt was designed and evaluated on the dataset, where it produced the best overall result.

\begin{thebibliography}{9}

\bibitem{chauhan2025mind}
Vashu Chauhan, Shobhit Raj, Shashank Mujumdar, Avirup Saha, and Anannay Jain.
\newblock \emph{Mind the Query: A Benchmark Dataset towards Text2Cypher Task}.
\newblock Proceedings of the 2025 Conference on Empirical Methods in Natural Language Processing: Industry Track, 2025.

\bibitem{hornsteiner2024realtime}
Markus Hornsteiner, Michael Kreussel, Christoph Steindl, Fabian Ebner, Philip Empl, and Stefan Sch{\"o}nig.
\newblock \emph{Real-Time Text-to-Cypher Query Generation with Large Language Models for Graph Databases}.
\newblock Future Internet, 16(12):438, 2024.

\bibitem{soleymani2025tasksplitting}
Saber Soleymani, Nathan Gravel, Krzysztof Kochut, and Natarajan Kannan.
\newblock \emph{Task Splitting and Prompt Engineering for Cypher Query Generation in Domain-Specific Knowledge Graphs}.
\newblock 2025.

\bibitem{mandilara2025decoding}
Ioanna Mandilara, Christina Maria Androna, Eleni Fotopoulou, Anastasios Zafeiropoulos, and Symeon Papavassiliou.
\newblock \emph{Decoding the Mystery: How Can LLMs Turn Text Into Cypher in Complex Knowledge Graphs?}
\newblock IEEE Access, 2025.

\bibitem{ozsoy2025text2cypher}
Makbule Gulcin Ozsoy et al.
\newblock \emph{Text2Cypher: Bridging Natural Language and Graph Databases}.
\newblock Proceedings of the ACL 2025 Workshop on GenAIK, 2025.

\end{thebibliography}

\end{document}


